\SourceLeaf{253}%
suppeditabunt determinationem coëfficientium \(\alpha\mathord{.}\ \beta,\ \gamma\mathord{.}\ \delta\). Erit scilicet
\[
\begin{aligned}
\alpha&=\BareRoot{\frac{E}{E'}}\mathord{.}\frac{\sin(\omega+\omega'-\psi)}{\sin\omega'},&\quad\beta&=\BareRoot{\frac{G}{E'}}\mathord{.}\frac{\sin(\omega'-\psi)}{\sin\omega'}\\
\gamma&=\BareRoot{\frac{E}{G'}}\mathord{.}\frac{\sin(\psi-\omega)}{\sin\omega'},&\delta&=\BareRoot{\frac{G}{G'}}\mathord{.}\frac{\sin\psi}{\sin\omega'}
\end{aligned}
\]
Adiungi debent aequationes
\[
\begin{aligned}
\cos\omega&=\frac{F}{\BareRoot{EG}},&\cos\omega'&=\frac{F'}{\BareRoot{E'G'}},\\
\sin\omega&=\BareRoot{\frac{EG-FF}{EG}},&\sin\omega'&=\BareRoot{\frac{E'G'-F'F'}{E'G'}}
\end{aligned}
\]
unde quatuor aequationes ita quoque exhiberi possunt
\[
\begin{aligned}
\alpha\BareRoot{(E'G'-F'F')}&=\BareRoot{EG'}\mathord{.}\sin(\omega+\omega'-\psi)\\
\beta\BareRoot{(E'G'-F'F')}&=\BareRoot{GG'}\mathord{.}\sin(\omega'-\psi)\\
\gamma\BareRoot{(E'G'-F'F')}&=\BareRoot{EE'}\mathord{.}\sin(\psi-\omega)\\
\delta\BareRoot{(E'G'-F'F')}&=\BareRoot{GE'}\mathord{.}\sin\psi
\end{aligned}
\]

Quum per substitutiones \(\mathrm{d}p'=\alpha\mathrm{d}p+\beta\mathrm{d}q\), \(\mathrm{d}q'=\gamma\mathrm{d}p+\delta\mathrm{d}q\) trinomium \(E'\mathrm{d}p'^2+2F'\mathrm{d}p'\mathord{.}\mathrm{d}q'+G'\mathrm{d}q'^2\) transire debeat in \(E\mathrm{d}p^2+2F\mathrm{d}p\mathord{.}\mathrm{d}q+G\mathrm{d}q^2\), facile obtinemus
\[
EG-FF=(E'G'-F'F')(\alpha\delta-\beta\gamma)^2
\]
et quum vice versa trinomium posterius rursus transire debeat in prius per substitutionem
\[
(\alpha\delta-\beta\gamma)\mathrm{d}p=\delta\mathrm{d}p'-\beta\mathrm{d}q',\qquad(\alpha\delta-\beta\gamma)\mathrm{d}q=-\gamma\mathrm{d}p'+\alpha\mathrm{d}q'
\]
invenimus
\[
\begin{aligned}
E\delta\delta-2F\gamma\delta+G\gamma\gamma&=\frac{EG-FF}{E'G'-F'F'}\mathord{.}E'\\
E\beta\delta-F(\alpha\delta+\beta\gamma)+G\alpha\gamma&=-\frac{EG-FF}{E'G'-F'F'}\mathord{.}F'\\
E\beta\beta-2F\alpha\beta+G\alpha\alpha&=\frac{EG-FF}{E'G'-F'F'}\mathord{.}G'
\end{aligned}
\]
\Article{22}
A disquisitione generali art. praec. descendimus ad applicationem latissime patentem, ubi, dum \(p\) et \(q\) etiam significatione generalissima accipiuntur, pro \(p',q'\) adoptamus quantitates in art. 15 per \(r,\varphi\) denotatas, quibus characteribus
